\IfFileExists{stacks-project.cls}{%
\documentclass{stacks-project}
}{%
\documentclass{amsart}
}

% For dealing with references we use the comment environment
\usepackage{verbatim}
\newenvironment{reference}{\comment}{\endcomment}
%\newenvironment{reference}{}{}
\newenvironment{slogan}{\comment}{\endcomment}
\newenvironment{history}{\comment}{\endcomment}

% For commutative diagrams we use Xy-pic
\usepackage[all]{xy}

% We use 2cell for 2-commutative diagrams.
\xyoption{2cell}
\UseAllTwocells

% We use multicol for the list of chapters between chapters
\usepackage{multicol}

% This is generally recommended for better output
\usepackage{lmodern}
\usepackage[T1]{fontenc}

% For cross-file-references
\usepackage{xr-hyper}

% Package for hypertext links:
\usepackage{hyperref}

% For any local file, say "hello.tex" you want to link to please
% use \externaldocument[hello-]{hello}
\externaldocument[introduction-]{introduction}
\externaldocument[conventions-]{conventions}
\externaldocument[sets-]{sets}
\externaldocument[categories-]{categories}
\externaldocument[topology-]{topology}
\externaldocument[sheaves-]{sheaves}
\externaldocument[sites-]{sites}
\externaldocument[stacks-]{stacks}
\externaldocument[fields-]{fields}
\externaldocument[algebra-]{algebra}
\externaldocument[brauer-]{brauer}
\externaldocument[homology-]{homology}
\externaldocument[derived-]{derived}
\externaldocument[simplicial-]{simplicial}
\externaldocument[more-algebra-]{more-algebra}
\externaldocument[smoothing-]{smoothing}
\externaldocument[modules-]{modules}
\externaldocument[sites-modules-]{sites-modules}
\externaldocument[injectives-]{injectives}
\externaldocument[cohomology-]{cohomology}
\externaldocument[sites-cohomology-]{sites-cohomology}
\externaldocument[dga-]{dga}
\externaldocument[dpa-]{dpa}
\externaldocument[sdga-]{sdga}
\externaldocument[hypercovering-]{hypercovering}
\externaldocument[schemes-]{schemes}
\externaldocument[constructions-]{constructions}
\externaldocument[properties-]{properties}
\externaldocument[morphisms-]{morphisms}
\externaldocument[coherent-]{coherent}
\externaldocument[divisors-]{divisors}
\externaldocument[limits-]{limits}
\externaldocument[varieties-]{varieties}
\externaldocument[topologies-]{topologies}
\externaldocument[descent-]{descent}
\externaldocument[perfect-]{perfect}
\externaldocument[more-morphisms-]{more-morphisms}
\externaldocument[flat-]{flat}
\externaldocument[groupoids-]{groupoids}
\externaldocument[more-groupoids-]{more-groupoids}
\externaldocument[etale-]{etale}
\externaldocument[chow-]{chow}
\externaldocument[intersection-]{intersection}
\externaldocument[pic-]{pic}
\externaldocument[weil-]{weil}
\externaldocument[adequate-]{adequate}
\externaldocument[dualizing-]{dualizing}
\externaldocument[duality-]{duality}
\externaldocument[discriminant-]{discriminant}
\externaldocument[derham-]{derham}
\externaldocument[local-cohomology-]{local-cohomology}
\externaldocument[algebraization-]{algebraization}
\externaldocument[curves-]{curves}
\externaldocument[resolve-]{resolve}
\externaldocument[models-]{models}
\externaldocument[functors-]{functors}
\externaldocument[equiv-]{equiv}
\externaldocument[pione-]{pione}
\externaldocument[etale-cohomology-]{etale-cohomology}
\externaldocument[proetale-]{proetale}
\externaldocument[relative-cycles-]{relative-cycles}
\externaldocument[more-etale-]{more-etale}
\externaldocument[trace-]{trace}
\externaldocument[crystalline-]{crystalline}
\externaldocument[spaces-]{spaces}
\externaldocument[spaces-properties-]{spaces-properties}
\externaldocument[spaces-morphisms-]{spaces-morphisms}
\externaldocument[decent-spaces-]{decent-spaces}
\externaldocument[spaces-cohomology-]{spaces-cohomology}
\externaldocument[spaces-limits-]{spaces-limits}
\externaldocument[spaces-divisors-]{spaces-divisors}
\externaldocument[spaces-over-fields-]{spaces-over-fields}
\externaldocument[spaces-topologies-]{spaces-topologies}
\externaldocument[spaces-descent-]{spaces-descent}
\externaldocument[spaces-perfect-]{spaces-perfect}
\externaldocument[spaces-more-morphisms-]{spaces-more-morphisms}
\externaldocument[spaces-flat-]{spaces-flat}
\externaldocument[spaces-groupoids-]{spaces-groupoids}
\externaldocument[spaces-more-groupoids-]{spaces-more-groupoids}
\externaldocument[bootstrap-]{bootstrap}
\externaldocument[spaces-pushouts-]{spaces-pushouts}
\externaldocument[spaces-chow-]{spaces-chow}
\externaldocument[groupoids-quotients-]{groupoids-quotients}
\externaldocument[spaces-more-cohomology-]{spaces-more-cohomology}
\externaldocument[spaces-simplicial-]{spaces-simplicial}
\externaldocument[spaces-duality-]{spaces-duality}
\externaldocument[formal-spaces-]{formal-spaces}
\externaldocument[restricted-]{restricted}
\externaldocument[spaces-resolve-]{spaces-resolve}
\externaldocument[formal-defos-]{formal-defos}
\externaldocument[defos-]{defos}
\externaldocument[cotangent-]{cotangent}
\externaldocument[examples-defos-]{examples-defos}
\externaldocument[algebraic-]{algebraic}
\externaldocument[examples-stacks-]{examples-stacks}
\externaldocument[stacks-sheaves-]{stacks-sheaves}
\externaldocument[criteria-]{criteria}
\externaldocument[artin-]{artin}
\externaldocument[quot-]{quot}
\externaldocument[stacks-properties-]{stacks-properties}
\externaldocument[stacks-morphisms-]{stacks-morphisms}
\externaldocument[stacks-limits-]{stacks-limits}
\externaldocument[stacks-cohomology-]{stacks-cohomology}
\externaldocument[stacks-perfect-]{stacks-perfect}
\externaldocument[stacks-introduction-]{stacks-introduction}
\externaldocument[stacks-more-morphisms-]{stacks-more-morphisms}
\externaldocument[stacks-geometry-]{stacks-geometry}
\externaldocument[moduli-]{moduli}
\externaldocument[moduli-curves-]{moduli-curves}
\externaldocument[examples-]{examples}
\externaldocument[exercises-]{exercises}
\externaldocument[guide-]{guide}
\externaldocument[desirables-]{desirables}
\externaldocument[coding-]{coding}
\externaldocument[obsolete-]{obsolete}
\externaldocument[fdl-]{fdl}
\externaldocument[index-]{index}

% Theorem environments.
%
\theoremstyle{plain}
\newtheorem{theorem}[subsection]{Theorem}
\newtheorem{proposition}[subsection]{Proposition}
\newtheorem{lemma}[subsection]{Lemma}

\theoremstyle{definition}
\newtheorem{definition}[subsection]{Definition}
\newtheorem{example}[subsection]{Example}
\newtheorem{exercise}[subsection]{Exercise}
\newtheorem{situation}[subsection]{Situation}

\theoremstyle{remark}
\newtheorem{remark}[subsection]{Remark}
\newtheorem{remarks}[subsection]{Remarks}

\numberwithin{equation}{subsection}

% Macros
%
\def\lim{\mathop{\mathrm{lim}}\nolimits}
\def\colim{\mathop{\mathrm{colim}}\nolimits}
\def\Spec{\mathop{\mathrm{Spec}}}
\def\Hom{\mathop{\mathrm{Hom}}\nolimits}
\def\Ext{\mathop{\mathrm{Ext}}\nolimits}
\def\SheafHom{\mathop{\mathcal{H}\!\mathit{om}}\nolimits}
\def\SheafExt{\mathop{\mathcal{E}\!\mathit{xt}}\nolimits}
\def\Sch{\mathit{Sch}}
\def\Mor{\mathop{\mathrm{Mor}}\nolimits}
\def\Ob{\mathop{\mathrm{Ob}}\nolimits}
\def\Sh{\mathop{\mathit{Sh}}\nolimits}
\def\NL{\mathop{N\!L}\nolimits}
\def\CH{\mathop{\mathrm{CH}}\nolimits}
\def\proetale{{pro\text{-}\acute{e}tale}}
\def\etale{{\acute{e}tale}}
\def\QCoh{\mathit{QCoh}}
\def\Ker{\mathop{\mathrm{Ker}}}
\def\Im{\mathop{\mathrm{Im}}}
\def\Coker{\mathop{\mathrm{Coker}}}
\def\Coim{\mathop{\mathrm{Coim}}}

% Boxtimes
%
\DeclareMathSymbol{\boxtimes}{\mathbin}{AMSa}{"02}

%
% Macros for moduli stacks/spaces
%
\def\QCohstack{\mathcal{QC}\!\mathit{oh}}
\def\Cohstack{\mathcal{C}\!\mathit{oh}}
\def\Spacesstack{\mathcal{S}\!\mathit{paces}}
\def\Quotfunctor{\mathrm{Quot}}
\def\Hilbfunctor{\mathrm{Hilb}}
\def\Curvesstack{\mathcal{C}\!\mathit{urves}}
\def\Polarizedstack{\mathcal{P}\!\mathit{olarized}}
\def\Complexesstack{\mathcal{C}\!\mathit{omplexes}}
% \Pic is the operator that assigns to X its picard group, usage \Pic(X)
% \Picardstack_{X/B} denotes the Picard stack of X over B
% \Picardfunctor_{X/B} denotes the Picard functor of X over B
\def\Pic{\mathop{\mathrm{Pic}}\nolimits}
\def\Picardstack{\mathcal{P}\!\mathit{ic}}
\def\Picardfunctor{\mathrm{Pic}}
\def\Deformationcategory{\mathcal{D}\!\mathit{ef}}
\usepackage{cleveref}
\crefname{originalreference}{original reference}{original references}
\Crefname{originalreference}{Original reference}{Original references}
\makeatletter
\newlabel{lemma-describe-sheaves-simplicial-site-site}{{09WF}{source}{}{sourcepointer-1}{}}
\newlabel{lemma-describe-sheaves-simplicial-site-site@cref}{{[originalreference][0][]09WF}{[1][0][]source}{}{}{}}
\makeatother
\crefname{originalreference}{original reference}{original references}
\Crefname{originalreference}{Original reference}{Original references}
\makeatletter
\newlabel{lemma-restriction-to-components-site}{{09WG}{source}{}{sourcepointer-2}{}}
\newlabel{lemma-restriction-to-components-site@cref}{{[originalreference][0][]09WG}{[1][0][]source}{}{}{}}
\makeatother
\crefname{originalreference}{original reference}{original references}
\Crefname{originalreference}{Original reference}{Original references}
\makeatletter
\newlabel{remark-augmentation-site}{{0D6Z}{source}{}{sourcepointer-3}{}}
\newlabel{remark-augmentation-site@cref}{{[originalreference][0][]0D6Z}{[1][0][]source}{}{}{}}
\makeatother
\crefname{originallemma}{original lemma}{original lemmas}
\Crefname{originallemma}{Original lemma}{Original lemmas}
\makeatletter
\newlabel{simplicial-lemma-augmentation-howto}{{018H}{source}{}{sourcepointer-4}{}}
\newlabel{simplicial-lemma-augmentation-howto@cref}{{[originallemma][0][]018H}{[1][0][]source}{}{}{}}
\makeatother
\crefname{originallemma}{original lemma}{original lemmas}
\Crefname{originallemma}{Original lemma}{Original lemmas}
\makeatletter
\newlabel{sites-lemma-composition-cocontinuous}{{03L5}{source}{}{sourcepointer-5}{}}
\newlabel{sites-lemma-composition-cocontinuous@cref}{{[originallemma][0][]03L5}{[1][0][]source}{}{}{}}
\makeatother
\crefname{originallemma}{original lemma}{original lemmas}
\Crefname{originallemma}{Original lemma}{Original lemmas}
\makeatletter
\newlabel{sites-lemma-composition-morphisms-sites}{{03CB}{source}{}{sourcepointer-6}{}}
\newlabel{sites-lemma-composition-morphisms-sites@cref}{{[originallemma][0][]03CB}{[1][0][]source}{}{}{}}
\makeatother
\crefname{originalreference}{original reference}{original references}
\Crefname{originalreference}{Original reference}{Original references}
\makeatletter
\newlabel{situation-simplicial-site}{{09WE}{source}{}{sourcepointer-7}{}}
\newlabel{situation-simplicial-site@cref}{{[originalreference][0][]09WE}{[1][0][]source}{}{}{}}
\makeatother
\title{An augmentation towards D induces morphisms of topoi a\_n and a on the total site with a\_* computed as an equalizer}
\author{The Stacks Project Authors}
\date{stacks-project:spaces-simplicial.tex:local-snapshot-20261012; GNU FDL 1.2 or later}
\begin{document}
\maketitle
\noindent\textit{Editorial note.} All delivered mathematical text is the author's own, copied from the pinned Stacks Project chapter file sources/stacks-spaces-simplicial/pinned-original-source.tex; no statement or proof step was rewritten, reordered or removed.\par\medskip
\noindent\textit{Editorial note.} The counted proof is the author's complete argument: item (3) holds by the definition of a\_* as the displayed equalizer (source lines 1204-1214), and the relation a g\_n = a\_n (line 1217, printed 'f\_n' in the source, a typo for a\_n) is the one-line componentwise check (a g\_n)\textasciicircum{}\{-1\}G = g\_n\textasciicircum{}\{-1\}(a\textasciicircum{}\{-1\}G) = (a\textasciicircum{}\{-1\}G)\_n = a\_n\textasciicircum{}\{-1\}G using the componentwise description of a\textasciicircum{}\{-1\} and part (1) of the restriction-to-components lemma; the author's spellings (including 'coosimplicial' and 'adjuntion') are preserved verbatim; this lemma supplies the equalizer description used by cards L46151 and the chapter's descent arguments.\par\medskip
\noindent\textit{Editorial note.} The label lemma-augmentation-site names the counted statement, and the proof environment immediately following it in the source is delivered verbatim as its proof.\par\medskip
\noindent\textit{Editorial note.} The counted statement is carried by the author's own lemma environment in the chapter source.\par\medskip
\noindent\textit{Allowed standard prerequisites (editorial).}\begin{itemize}
\item The augmentation datum of the simplicial site. Same chapter, Remark \textbackslash{}label\{remark-augmentation-site\} and Situation \textbackslash{}ref\{situation-simplicial-site\}. Lines 1142-1145: the functors u\_0, u\_varphi with the relations u\_varphi u\_0-style compositions of the datum.
\item Sites, composition of morphisms of sites (case A). Sites chapter, \textbackslash{}label\{sites-lemma-composition-morphisms-sites\}. Lines 1163-1169: u\_n = u\_varphi u\_0 defines the morphism of sites a\_n and gives a\_m f\_varphi = a\_n.
\item Sites, composition of cocontinuous functors (case B). Sites chapter, \textbackslash{}label\{sites-lemma-composition-cocontinuous\}. Lines 1171-1177: the cocontinuous u\_n defines a\_n : Sh(C\_n) -> Sh(D) with the same compatibilities.
\item Sheaves on the total site described via components and transitions. Same chapter, \textbackslash{}label\{lemma-describe-sheaves-simplicial-site-site\}. Lines 1189-1195: the componentwise definition of a\textasciicircum{}\{-1\}G with transition maps the identifications f\_varphi\textasciicircum{}\{-1\} a\_m\textasciicircum{}\{-1\}G = a\_n\textasciicircum{}\{-1\}G.
\item The cosimplicial augmentation lemma for maps into cosimplicial objects. Simplicial chapter, \textbackslash{}label\{simplicial-lemma-augmentation-howto\} (dual). Lines 1222-1226: maps G\_bullet -> a\_\{bullet,*\}F of cosimplicial sheaves are the same as maps G -> a\_*F.
\item The component morphisms g\_n of the total site. Same chapter, \textbackslash{}label\{lemma-restriction-to-components-site\} (card L46149 of this batch). Lines 1152-1155: item (2) of the statement references these g\_n; a g\_n = a\_n is checked from the componentwise description.
\end{itemize}\medskip
\ifdefined\hypertarget\hypertarget{sourcepointer-1}{}\fi
\noindent\textit{Source reference.} Original reference 09WF: sources/stacks-spaces-simplicial/pinned-original-source.tex, source lines 953--953. Prerequisite pointer invoked at proof line(s) 1193; it names a standard fact of the same chapter that the counted argument uses, and every step of the counted proof built on it is retained verbatim in the delivered body.\par\medskip
\ifdefined\hypertarget\hypertarget{sourcepointer-2}{}\fi
\noindent\textit{Source reference.} Original reference 09WG: sources/stacks-spaces-simplicial/pinned-original-source.tex, source lines 970--970. Prerequisite pointer invoked at proof line(s) 1154; it names a standard fact of the same chapter that the counted argument uses, and every step of the counted proof built on it is retained verbatim in the delivered body.\par\medskip
\ifdefined\hypertarget\hypertarget{sourcepointer-3}{}\fi
\noindent\textit{Source reference.} Original reference 0D6Z: sources/stacks-spaces-simplicial/pinned-original-source.tex, source lines 1125--1125. Prerequisite pointer invoked at proof line(s) 1144; it names a standard fact of the same chapter that the counted argument uses, and every step of the counted proof built on it is retained verbatim in the delivered body.\par\medskip
\ifdefined\hypertarget\hypertarget{sourcepointer-4}{}\fi
\noindent\textit{Source reference.} Original lemma 018H: sources/stacks-chow/tags.txt, source lines 1616--1616. Prerequisite pointer invoked at proof line(s) 1225; it names a standard fact of a sibling chapter that the counted argument uses, and every step of the counted proof built on it is retained verbatim in the delivered body.\par\medskip
\ifdefined\hypertarget\hypertarget{sourcepointer-5}{}\fi
\noindent\textit{Source reference.} Original lemma 03L5: sources/stacks-chow/tags.txt, source lines 4701--4701. Prerequisite pointer invoked at proof line(s) 1176; it names a standard fact of a sibling chapter that the counted argument uses, and every step of the counted proof built on it is retained verbatim in the delivered body.\par\medskip
\ifdefined\hypertarget\hypertarget{sourcepointer-6}{}\fi
\noindent\textit{Source reference.} Original lemma 03CB: sources/stacks-chow/tags.txt, source lines 4378--4378. Prerequisite pointer invoked at proof line(s) 1167; it names a standard fact of a sibling chapter that the counted argument uses, and every step of the counted proof built on it is retained verbatim in the delivered body.\par\medskip
\ifdefined\hypertarget\hypertarget{sourcepointer-7}{}\fi
\noindent\textit{Source reference.} Original reference 09WE: sources/stacks-spaces-simplicial/pinned-original-source.tex, source lines 888--888. Prerequisite pointer invoked at proof line(s) 1142; it names a standard fact of the same chapter that the counted argument uses, and every step of the counted proof built on it is retained verbatim in the delivered body.\par\medskip
\begin{lemma}
\label{lemma-augmentation-site}
In Situation \ref{situation-simplicial-site} let $a_0$ be an
augmentation towards a site $\mathcal{D}$ as in
Remark \ref{remark-augmentation-site}. Then $a_0$ induces
\begin{enumerate}
\item a morphism of topoi $a_n : \Sh(\mathcal{C}_n) \to \Sh(\mathcal{D})$
for all $n \geq 0$,
\item a morphism of topoi $a : \Sh(\mathcal{C}_{total}) \to \Sh(\mathcal{D})$
\end{enumerate}
such that
\begin{enumerate}
\item for all $\varphi : [m] \to [n]$ we have $a_m \circ f_\varphi = a_n$,
\item if $g_n : \Sh(\mathcal{C}_n) \to \Sh(\mathcal{C}_{total})$
is as in Lemma \ref{lemma-restriction-to-components-site}, then
$a \circ g_n = a_n$, and
\item $a_*\mathcal{F}$ for $\mathcal{F} \in \Sh(\mathcal{C}_{total})$
is the equalizer of the two maps
$a_{0, *}\mathcal{F}_0 \to a_{1, *}\mathcal{F}_1$.
\end{enumerate}
\end{lemma}

\begin{proof}
Case A. Let $u_n : \mathcal{D} \to \mathcal{C}_n$ be the common
value of the functors $u_\varphi \circ u_0$ for $\varphi : [0] \to [n]$.
Then $u_n$ corresponds to a morphism of sites
$a_n : \mathcal{C}_n \to \mathcal{D}$, see
Sites, Lemma \ref{sites-lemma-composition-morphisms-sites}.
The same lemma shows that for all $\varphi : [m] \to [n]$ we have
$a_m \circ f_\varphi = a_n$.

\medskip\noindent
Case B. Let $u_n : \mathcal{C}_n \to \mathcal{D}$ be the common
value of the functors $u_0 \circ u_\varphi$ for $\varphi : [0] \to [n]$.
Then $u_n$ is cocontinuous and hence defines a morphism of topoi
$a_n : \Sh(\mathcal{C}_n) \to \Sh(\mathcal{D)}$, see
Sites, Lemma \ref{sites-lemma-composition-cocontinuous}.
The same lemma shows that for all $\varphi : [m] \to [n]$ we have
$a_m \circ f_\varphi = a_n$.

\medskip\noindent
Consider the functor $a^{-1} : \Sh(\mathcal{D}) \to \Sh(\mathcal{C}_{total})$
which to a sheaf of sets $\mathcal{G}$ associates the sheaf
$\mathcal{F} = a^{-1}\mathcal{G}$ whose components are $a_n^{-1}\mathcal{G}$
and whose transition maps $\mathcal{F}(\varphi)$ are the identifications
$$
f_\varphi^{-1}\mathcal{F}_m =
f_\varphi^{-1} a_m^{-1}\mathcal{G} =
a_n^{-1}\mathcal{G} =
\mathcal{F}_n
$$
for $\varphi : [m] \to [n]$, see the description of
$\Sh(\mathcal{C}_{total})$ in
Lemma \ref{lemma-describe-sheaves-simplicial-site-site}.
Since the functors $a_n^{-1}$ are exact, $a^{-1}$ is an exact functor.
Finally, for $a_* : \Sh(\mathcal{C}_{total}) \to \Sh(\mathcal{D})$
we take the functor which to a sheaf $\mathcal{F}$ on $\Sh(\mathcal{D})$
associates
$$
\xymatrix{
a_*\mathcal{F} \ar@{=}[r] &
\text{Equalizer}(a_{0, *}\mathcal{F}_0
\ar@<1ex>[r] \ar@<-1ex>[r] &
a_{1, *}\mathcal{F}_1)
}
$$
Here the two maps come from the two maps $\varphi : [0] \to [1]$
via
$$
a_{0, *}\mathcal{F}_0 \to
a_{0, *}f_{\varphi, *} f_\varphi^{-1}\mathcal{F}_0
\xrightarrow{\mathcal{F}(\varphi)}
a_{0, *}f_{\varphi, *} \mathcal{F}_1 = a_{1, *}\mathcal{F}_1
$$
where the first arrow comes from $1 \to f_{\varphi, *} f_\varphi^{-1}$.
Let $\mathcal{G}_\bullet$ denote the constant coosimplicial sheaf
with value $\mathcal{G}$ and let $a_{\bullet, *}\mathcal{F}$
denote the cosimplicial sheaf having $a_{n, *}\mathcal{F}_n$ in degree $n$.
By the usual adjuntion for the morphisms of topoi $a_n$ we see that
a map $a^{-1}\mathcal{G} \to \mathcal{F}$
is the same thing as a map
$$
\mathcal{G}_\bullet \longrightarrow a_{\bullet, *}\mathcal{F}
$$
of cosimplicial sheaves.
By the dual to Simplicial, Lemma \ref{simplicial-lemma-augmentation-howto}
this is the same thing as a map $\mathcal{G} \to a_*\mathcal{F}$.
Thus $a^{-1}$ and $a_*$ are adjoint functors and we obtain
our morphism of topoi $a$\footnote{In case B the morphism $a$
corresponds to the cocontinuous functor
$\mathcal{C}_{total} \to \mathcal{D}$ sending
$U$ in $\mathcal{C}_n$ to $u_n(U)$.}. The equalities
$a \circ g_n = f_n$ follow immediately from the definitions.
\end{proof}

\end{document}
